\section{Multi-Criteria Analysis}\label{sec:multi-criteria-analysis}

\subsection{Criteria}\label{subsec:criteria}
A set of criteria was established to evaluate the back-end frameworks.
The reasoning behind each criterion, as well as their assigned weights, are explained below.

\subsubsection{Communication Protocols}\label{subsubsec:communication-protocols}
The to be built software solution requires two types of communication.

The back-end application will provide an API that follows the REST protocol standards.
Since the front-end application needs to communicate with the back-end to perform CRUD operations, it must be able to construct and send \emph{RESTful API calls}.

Additionally, the software solution will include chat functionality.
For real-time communication between users, some sort of socket connection is required.
This means the application must also be able to establish and maintain a \emph{WebSocket connection}.

\paragraph{Method of evaluation}
To score this criterion, packages or built-in solutions will be looked up for both RESTful API calls and WebSocket connections.
In case any of these functionalities are not supported, the framework will be disqualified.

\paragraph{Weight}
This criterion is not assigned any weight.

If the framework does not give the option to perform both acts of communication, the functionality of the software solution cannot be realized, meaning the application cannot be built using that framework.

\subsubsection{Community}\label{subsubsec:community}
The larger the community, the more online resources and support will be available for learning and troubleshooting the framework.
To assess the size of the community, several metrics will be considered:

\begin{itemize}
	\item The \emph{number of stars} on the official framework GitHub repository.
	      While not a direct measure of usage, it gives a rough indication of how many developers are interested in or actively using the framework.
	      The more developers using the framework, the more questions that will have been asked and answered online.
	\item The number of \emph{Discord users} in the official or related Discord servers.
	      A larger user count means more potential help when asking questions.
	\item The number of \emph{examples} online.
	      Examples of implementations, tutorials, and blog posts can be very helpful when learning a new framework.
	      If the framework is unopinionated, a big variety of examples can help in deciding how to structure the application.
\end{itemize}

\paragraph{Method of evaluation}
To score this criterion, the stars from the official GitHub repository will be counted, the official or related Discord servers will be joined to view the member count, and a search on GitHub will be performed to find the number of repositories using the framework.

\paragraph{Weight}
This criterion is assigned a weight of \emph{20}, derived from the above sub-criteria, which are \emph{5}, \emph{5}, and \emph{10}.

Both the number of stars and the number of Discord users are not very important, because an application can be built without community support as well.
It would just be more difficult.

The number of available examples holds slightly more value, as it provides real-world insights into project structure and professional best practices.
That said, while they can speed up the learning process, an application can still be built without them.

For these reasons, the community criterion is not assigned a very high weight.

\subsubsection{Developer Experience}\label{subsubsec:developer-experience}
A good developer experience is important for productivity and ease of use.
Even if a framework is powerful and flexible, if it is difficult to use, it will slow down development and may frustrate the developer.
It should be easy to work with, so developers can focus on building the application rather than struggling with the tooling.
The following aspects will be considered to determine good developer experience:

\begin{itemize}
	\item The framework and the programming language it uses should not have a high \emph{learning curve}.
	      If a framework differs significantly from other frameworks, one which the developer is not yet familiar with, building an application within a short time becomes challenging.
	\item There should be built-in support for \emph{debugging} the application during runtime using a Jetbrains IDE.
	      Without proper debugging tools, identifying and fixing issues can be time-consuming and difficult.
	\item High quality \emph{documentation} should be readily available.
	      Poor documentation forces developers to rely on user discussions and examples, which may not always be accurate or up-to-date.
	      Good documentation should include clear explanations, code examples, and guides on best practices.
\end{itemize}

\paragraph{Method of evaluation}
To score this criterion, the existence of a debugging solution as well as the presence of proper documentation will be checked.
To validate that the documentation is proper, the website must have some sort of installation guide, and documentation on specific parts of the framework, such as routing and middleware.

The other part, the learning curve, will be determined subjectively based on the developer's prior experience with similar frameworks and programming languages.
The proficiency scale is as follows:

\begin{itemize}
	\item \textbf{0.00} - No prior experience with the framework or programming language.
	\item \textbf{0.25} - Limited experience with the framework or programming language.
	      The developer has worked with it on small projects or has only read about it.
	\item \textbf{0.50} - Moderate experience with the framework or programming language.
	      The developer has worked with it on a few projects and has a basic understanding of its features and best practices.
	\item \textbf{0.75} - Good experience with the framework or programming language.
	      The developer has worked with it on several projects and is familiar with its features and best practices.
	\item \textbf{1.00} - Extensive experience with the framework or programming language.
	      The developer has worked with it on many projects and is considered an expert in its use.
\end{itemize}

\paragraph{Weight}
This criterion is assigned a weight of \emph{45}, derived from the above sub-criteria, which are \emph{20}, \emph{10}, and \emph{15}.

The learning curve is quite important.
Time is the limiting factor in this project, and both prototypes and the final application need to be built within a few weeks.
If a framework takes too long to learn, it will not be feasible to use it.

Debugging is also important, but even without it, development is still very doable.
Most programming languages have some sort of logging functionality, which can be used to identify issues instead.

Documentation is very important, especially when working with a framework for the first time.
While a large and proficient community can provide answers, it is better to have standardized information readily available rather than having to rely solely on other users.

For these reasons, the developer experience criterion is assigned a relatively high weight.

\subsubsection{Ecosystem}\label{subsubsec:ecosystem}
The ecosystem comprises all available resources to reduce the workload of developers.
The more resources are available, the more developers can focus on implementing business logic rather than reinventing the wheel.
The following aspects will be considered to determine the quality of the ecosystem:

\begin{itemize}
	\item There should be official or well-maintained community \emph{plugins} for Jetbrains IDEs, offering features such as code completion, tooltip documentation, and error checking.
	      These plugins can significantly enhance productivity and reduce the likelihood of errors.
	\item A wide variety of \emph{libraries} should be available.
	      These libraries can provide pre-built functionality, saving developers time and effort.
	      The higher the number of available libraries, the bigger the chance that a library exists for a specific need.
\end{itemize}

\paragraph{Method of evaluation}
To score this criterion, the availability of IDE plugin support will be checked, as well as the number of existing libraries.

\paragraph{Weight}
This criterion is assigned a weight of \emph{40}, derived from the above sub-criteria, which are \emph{25} and \emph{15}.

The availability of IDE plugins is very important.
Having tools that streamline development makes a big difference.
If these tools are not available, development will be slower and more error-prone.

The number of libraries is also important, but to a lesser extent.
While they can speed up development, developers can create their own libraries if needed.
This however, takes more time, making a rich ecosystem of ready-to-use libraries preferable.

For these reasons, the ecosystem criterion is assigned a high weight.

\subsubsection{Quality Assurance}\label{subsubsec:quality-assurance}
The app should be of good quality when released to the public, ensuring both solid security and proper functionality.
This criterion encapsulates a part of the advanced track chosen for back-end development (\emph{Quality Management and Deployment}).
To help ensure the quality of the application, the following aspects will be considered:

\begin{itemize}
	\item The framework should provide some form of (automated) \emph{testing}.
	      If this is not done, new updates may introduce bugs or break existing functionality, which could confuse or frustrate the users.
	\item The framework should provide some form of \emph{static code analysis} tooling.
	      This ensures the code is of good quality and keeps consistency between code that is written by different developers.
	\item The framework should receive regular \emph{updates}.
	      If bugs or security flaws are discovered, they should be fixed timely.
	      Without frequent updates, users are left at risk, while developers have no control of critical issues.
\end{itemize}

\paragraph{Method of evaluation}
To score this criterion, the existence of unit and integration testing solutions, and linting tools, will be checked.
In addition, it will be verified that the framework has received an update within the last six months.

\paragraph{Weight}
This criterion is assigned a weight of \emph{25}, derived from the above sub-criteria, which are \emph{15}, \emph{5}, and \emph{5}.

Testing the application is rather important, since it directly affects the quality of the deployed application from a user's perspective.

Static code analysis is not very important.
While it helps maintain cleaner and more consistent code, it does not directly affect the quality of the deployed application from a user's perspective.
Since this project is small-scale, the codebase will be manageable without it.

Updates are less important as long as the application remains secure and free of major bugs.
However, the framework should still be maintained to ensure long-term stability.

For these reasons, the quality assurance criterion is assigned a moderate weight.

\subsection{Frameworks}\label{subsec:frameworks}
The following back-end frameworks will be evaluated based on the criteria established in \cref{subsec:criteria}.

\subsubsection{Echo}\label{subsubsec:frameworks-echo}
Echo is a minimalist Go web framework by Labstack.
It is lightweight, fast, and supports middleware for authentication, logging, and request validation.
It includes built-in support for JWT, CORS, and WebSockets.
Echo was chosen for its simplicity, performance, and extendability. \cite{bib:echo-framework}

\subsubsection{Spring Boot}\label{subsubsec:frameworks-spring-boot}
Spring Boot is an open source Java-based framework developed by VMware \cite{bib:springboot-creator}, built on top of the Spring Framework.
It simplifies the configuration and deployment of production-ready applications by providing auto-configuration, embedded servers, and a rich ecosystem of starter dependencies. \cite{bib:spring-boot-framework}

Spring Boot was chosen because it is one of the most widely used back-end frameworks in enterprise development, and because it is built on Java, a language both team members have prior experience with.

\subsubsection{NestJS}\label{subsubsec:frameworks-nestjs}
NestJS is a progressive Node.js framework for building efficient and scalable server-side applications using TypeScript.
It is heavily inspired by Angular and uses decorators, modules, and dependency injection to provide a structured and opinionated architecture. \cite{bib:nestjs-framework}

NestJS was chosen because it is built on TypeScript and JavaScript, which are already familiar to the team from front-end development, making knowledge transfer straightforward.

\subsection{Evaluation}\label{subsec:evaluation}
The frameworks are evaluated based on the criteria established in \cref{subsec:criteria}.

\subsubsection{Echo}\label{subsubsec:evaluation-echo}

\paragraph{Communication Protocols}
Golang has a built-in HTTP client that can perform REST calls. \cite{bib:golang-http-client}
For WebSockets, the \emph{gorilla/websocket} package provides a fast, well-tested WebSocket implementation for Go. \cite{bib:gorilla-websocket}

Since this criterion is not assigned any weight, Echo meets the requirements and passes the criterion.

\paragraph{Community}
The official Echo GitHub repository has over 32.200 stars \cite{bib:echo-github}, there is no dedicated Echo Discord server but the Golang Discord server has over 42.000 members \cite{bib:golang-discord}, and there are at least 6.500 results on GitHub for Echo projects \cite{bib:echo-github-search}.

The scores for this criterion are shown in \cref{tab:multi-criteria-analysis-echo-community}.

\begin{htable}[
		caption = {Echo community evaluation},
		label = {tab:multi-criteria-analysis-echo-community},
	]{
		colspec={lllX},
	}
	Sub-criterion & Weight & Function & Normalized score \\
	Number of stars & 5 & $\bar{x} = \frac{N}{N_{\max}}$ & $\bar{x} = \frac{32200}{32200} = 1$ \\
	Number of Discord users & 5 & $\bar{x} = \frac{N}{N_{\max}}$ & $\bar{x} = \frac{42000}{44000} = 0.95$ \\
	Number of examples & 10 & $\bar{x} = \frac{N}{N_{\max}}$ & $\bar{x} = \frac{6500}{508000} = 0.01$ \\
\end{htable}

\[
	S_{\text{community}} =
	\frac{
		5 \cdot 1 +
		5 \cdot 0.95 +
		10 \cdot 0.01
	}{
		5 + 5 + 10
	} = 0.55
\]

The normalized total score for this criterion is \emph{0.55}.

\paragraph{Developer Experience}
The proficiency ratings are as follows:
\begin{itemize}
	\item \textbf{Moustafa (0.60)} - Has built a REST API using Golang in his free time before, and feels comfortable with the language and its patterns.
	\item \textbf{William (0.05)} - Has not got any experience with Golang or Echo, but did review a REST API built with it, which gave him some idea of how the framework works.
\end{itemize}

Debugging Go applications can be done inside GoLand using JetBrains' built-in debugger. \cite{bib:goland-ide}
The documentation for Echo contains an installation guide and covers routing, middleware, and request handling, and the Go standard library has its own comprehensive documentation website as well. \cite{bib:echo-documentation}

The scores for this criterion are shown in \cref{tab:multi-criteria-analysis-echo-developer-experience}.

\begin{htable}[
		caption = {Echo developer experience evaluation},
		label = {tab:multi-criteria-analysis-echo-developer-experience},
	]{
		colspec={lllX},
	}
	Sub-criterion & Weight & Function & Normalized score \\
	Learning curve & 20 & $\bar{x} = \frac{x_1 + x_2}{2}$ & $\bar{x} = \frac{0.60 + 0.05}{2} = 0.33$ \\
	Debugging support & 10 & Binary (0 or 1) & 1 \\
	Documentation quality & 15 & Binary (0 or 1) & 1 \\
\end{htable}

\[
	S_{\text{developer experience}} =
	\frac{
		20 \cdot 0.33 +
		10 \cdot 1 +
		15 \cdot 1
	}{
		20 + 10 + 15
	} = 0.75
\]

The normalized total score for this criterion is \emph{0.75}.

\paragraph{Ecosystem}
JetBrains provides GoLand, a dedicated IDE for the Go programming language, which supports plugins for linting, testing, and other development tools. \cite{bib:goland-ide}
The Go package repository at \texttt{pkg.go.dev} hosts a large number of packages compatible with Echo applications.

The scores for this criterion are shown in \cref{tab:multi-criteria-analysis-echo-ecosystem}.

\begin{htable}[
		caption = {Echo ecosystem evaluation},
		label = {tab:multi-criteria-analysis-echo-ecosystem},
	]{
		colspec={lllX},
	}
	Sub-criterion & Weight & Function & Normalized score \\
	IDE plugin support & 25 & Binary (0 or 1) & 1 \\
	Number of libraries & 15 & $\bar{x} = \frac{N}{N_{\max}}$ & $\bar{x} = \frac{0}{2000} = 0$ \\
\end{htable}

\[
	S_{\text{ecosystem}} =
	\frac{
		25 \cdot 1 +
		15 \cdot 0
	}{
		25 + 15
	} = 0.63
\]

The normalized total score for this criterion is \emph{0.63}.

\paragraph{Quality Assurance}
Echo makes use of the Golang testing platform, which is built into the standard library and supports unit and integration testing. \cite{bib:golang-testing}
The Go ecosystem provides linting tools such as \emph{golangci-lint}, which bundles multiple linters and integrates with major JetBrains IDEs. \cite{bib:golangci-lint}
The last minor release of the Echo framework was the last six months. \cite{bib:echo-github}

The scores for this criterion are shown in \cref{tab:multi-criteria-analysis-echo-quality-assurance}.

\begin{htable}[
		caption = {Echo quality assurance evaluation},
		label = {tab:multi-criteria-analysis-echo-quality-assurance},
	]{
		colspec={lllX},
	}
	Sub-criterion & Weight & Function & Normalized score \\
	Testing & 15 & Binary (0 or 1) & 1 \\
	Static code analysis & 5 & Binary (0 or 1) & 1 \\
	Regular updates & 5 & Binary (0 or 1) & 1 \\
\end{htable}

\[
	S_{\text{quality assurance}} =
	\frac{
		15 \cdot 1 +
		5 \cdot 1 +
		5 \cdot 1
	}{
		15 + 5 + 5
	} = 1.00
\]

The normalized total score for this criterion is \emph{1.00}.

\subsubsection{Spring Boot}\label{subsubsec:evaluation-spring-boot}

\paragraph{Communication Protocols}
Spring Boot provides a built-in \texttt{RestTemplate} and the newer \texttt{WebClient} for performing RESTful API calls. \cite{bib:spring-boot-rest}
For WebSocket connections, Spring Boot includes the \emph{spring-websocket} module, which supports full-duplex communication using the STOMP messaging protocol. \cite{bib:spring-boot-websocket}

Since this criterion is not assigned any weight, Spring Boot meets the requirements and passes the criterion.

\paragraph{Community}
The official Spring Boot GitHub repository has over 80.000 stars \cite{bib:spring-boot-github}. Although there is no official Spring Boot Discord server, there are other communities focused on Java in general. Additionally, a search on GitHub returns over 400,000 Spring Boot projects \cite{bib:spring-boot-github-search}.

The scores for this criterion are shown in \cref{tab:multi-criteria-analysis-spring-boot-community}.

\begin{htable}[
		caption = {Spring Boot community evaluation},
		label = {tab:multi-criteria-analysis-spring-boot-community},
	]{
		colspec={lllX},
	}
	Sub-criterion & Weight & Function & Normalized score \\
	Number of stars & 5 & $\bar{x} = \frac{N}{N_{\max}}$ & $\bar{x} = \frac{80000}{80000} = 1$ \\
	Number of Discord users & 5 & $\bar{x} = \frac{N}{N_{\max}}$ & $\bar{x} = \frac{0}{0} = 0.0$ \\
	Number of examples & 10 & $\bar{x} = \frac{N}{N_{\max}}$ & $\bar{x} = \frac{400000}{508000} = 0.79$ \\
\end{htable}

\[
	S_{\text{community}} =
	\frac{
		5 \cdot 1 +
		5 \cdot 0.18 +
		10 \cdot 0.79
	}{
		5 + 5 + 10
	} = 0.69
\]

The normalized total score for this criterion is \emph{0.69}.

\paragraph{Developer Experience}
The proficiency ratings are as follows:
\begin{itemize}
	\item \textbf{Moustafa (0.25)} - Has some experience with Java from school, but has not used Spring Boot or the Spring ecosystem before.
	\item \textbf{William (0.25)} - Has limited experience with Java and has not used Spring Boot, though the general structure is somewhat familiar from other frameworks.
\end{itemize}

Debugging Spring Boot applications can be done inside IntelliJ IDEA, which provides first-class support for the Spring framework including run configurations and bean inspection. \cite{bib:intellij-spring}
The Spring Boot documentation is extensive and covers installation, configuration, routing, security, and data access. \cite{bib:spring-boot-documentation}

The scores for this criterion are shown in \cref{tab:multi-criteria-analysis-spring-boot-developer-experience}.

\begin{htable}[
		caption = {Spring Boot developer experience evaluation},
		label = {tab:multi-criteria-analysis-spring-boot-developer-experience},
	]{
		colspec={lllX},
	}
	Sub-criterion & Weight & Function & Normalized score \\
	Learning curve & 20 & $\bar{x} = \frac{x_1 + x_2}{2}$ & $\bar{x} = \frac{0.25 + 0.25}{2} = 0.25$ \\
	Debugging support & 10 & Binary (0 or 1) & 1 \\
	Documentation quality & 15 & Binary (0 or 1) & 1 \\
\end{htable}

\[
	S_{\text{developer experience}} =
	\frac{
		20 \cdot 0.25 +
		10 \cdot 1 +
		15 \cdot 1
	}{
		20 + 10 + 15
	} = 0.67
\]

The normalized total score for this criterion is \emph{0.67}.

\paragraph{Ecosystem}
IntelliJ IDEA Ultimate provides dedicated Spring Boot support including Spring-specific tooling, run configurations, and bean diagrams. \cite{bib:intellij-spring}
The Maven Central registry contains over 686.000 packages, a large portion of which are compatible with Spring Boot applications. \cite{bib:maven-central}

The scores for this criterion are shown in \cref{tab:multi-criteria-analysis-spring-boot-ecosystem}.

\begin{htable}[
		caption = {Spring Boot ecosystem evaluation},
		label = {tab:multi-criteria-analysis-spring-boot-ecosystem},
	]{
		colspec={lllX},
	}
	Sub-criterion & Weight & Function & Normalized score \\
	IDE plugin support & 25 & Binary (0 or 1) & 1 \\
	Number of libraries & 15 & $\bar{x} = \frac{N}{N_{\max}}$ & $\bar{x} = \frac{686}{2000} = 0.34$ \\
\end{htable}

\[
	S_{\text{ecosystem}} =
	\frac{
		25 \cdot 1 +
		15 \cdot 0.34
	}{
		25 + 15
	} = 0.75
\]

The normalized total score for this criterion is \emph{0.75}.

\paragraph{Quality Assurance}
Spring Boot integrates with JUnit and provides the \texttt{@SpringBootTest} annotation for integration testing, as well as MockMvc for testing web layers in isolation. \cite{bib:spring-boot-testing}
The Java ecosystem offers mature static analysis tools such as Checkstyle, PMD, and SpotBugs, all of which integrate well with JetBrains IntelliJ IDEA. \cite{bib:java-linting}
The last minor release of Spring Boot was within the last six months. \cite{bib:spring-boot-github}

The scores for this criterion are shown in \cref{tab:multi-criteria-analysis-spring-boot-quality-assurance}.

\begin{htable}[
		caption = {Spring Boot quality assurance evaluation},
		label = {tab:multi-criteria-analysis-spring-boot-quality-assurance},
	]{
		colspec={lllX},
	}
	Sub-criterion & Weight & Function & Normalized score \\
	Testing & 15 & Binary (0 or 1) & 1 \\
	Static code analysis & 5 & Binary (0 or 1) & 1 \\
	Regular updates & 5 & Binary (0 or 1) & 1 \\
\end{htable}

\[
	S_{\text{quality assurance}} =
	\frac{
		15 \cdot 1 +
		5 \cdot 1 +
		5 \cdot 1
	}{
		15 + 5 + 5
	} = 1.00
\]

The normalized total score for this criterion is \emph{1.00}.

\subsubsection{NestJS}\label{subsubsec:evaluation-nestjs}

\paragraph{Communication Protocols}
NestJS runs on Node.js, so either the Fetch API or \emph{Axios} can be used to perform RESTful API calls. \cite{bib:axios}
For WebSocket connections, NestJS has a built-in WebSockets module that supports both native WebSockets and Socket.IO. \cite{bib:nestjs-websockets}

Since this criterion is not assigned any weight, NestJS meets the requirements and passes the criterion.

\paragraph{Community}
The official NestJS GitHub repository has over 75.000 stars \cite{bib:nestjs-github}, the official NestJS Discord server has over 51.000 members \cite{bib:nestjs-discord}, and there are at least 80.000 results on GitHub for NestJS projects \cite{bib:nestjs-github-search}.

The scores for this criterion are shown in \cref{tab:multi-criteria-analysis-nestjs-community}.

\begin{htable}[
		caption = {NestJS community evaluation},
		label = {tab:multi-criteria-analysis-nestjs-community},
	]{
		colspec={lllX},
	}
	Sub-criterion & Weight & Function & Normalized score \\
	Number of stars & 5 & $\bar{x} = \frac{N}{N_{\max}}$ & $\bar{x} = \frac{75000}{75000} = 0.89$ \\
	Number of Discord users & 5 & $\bar{x} = \frac{N}{N_{\max}}$ & $\bar{x} = \frac{51000}{51000} = 1$ \\
	Number of examples & 10 & $\bar{x} = \frac{N}{N_{\max}}$ & $\bar{x} = \frac{80000}{508000} = 0.16$ \\
\end{htable}

\[
	S_{\text{community}} =
	\frac{
		5 \cdot 0.89 +
		5 \cdot 0.68 +
		10 \cdot 0.16
	}{
		5 + 5 + 10
	} = 0.47
\]

The normalized total score for this criterion is \emph{0.47}.

\paragraph{Developer Experience}
The proficiency ratings are as follows:
\begin{itemize}
	\item \textbf{Moustafa (0.25)} - Has minimal experience with JavaScript frameworks such as VueJS and Svelte, but has not used NestJS or TypeScript extensively.
	\item \textbf{William (0.50)} - Has experience with several JavaScript and TypeScript frameworks and feels that the Angular-inspired structure of NestJS is relatively familiar.
\end{itemize}

Debugging NestJS applications can be done inside WebStorm, which provides full TypeScript and Node.js debugging support. \cite{bib:webstorm-ide}
The NestJS documentation contains an installation guide and covers modules, controllers, providers, middleware, and WebSockets in detail. \cite{bib:nestjs-documentation}

The scores for this criterion are shown in \cref{tab:multi-criteria-analysis-nestjs-developer-experience}.

\begin{htable}[
		caption = {NestJS developer experience evaluation},
		label = {tab:multi-criteria-analysis-nestjs-developer-experience},
	]{
		colspec={lllX},
	}
	Sub-criterion & Weight & Function & Normalized score \\
	Learning curve & 20 & $\bar{x} = \frac{x_1 + x_2}{2}$ & $\bar{x} = \frac{0.25 + 0.50}{2} = 0.38$ \\
	Debugging support & 10 & Binary (0 or 1) & 1 \\
	Documentation quality & 15 & Binary (0 or 1) & 1 \\
\end{htable}

\[
	S_{\text{developer experience}} =
	\frac{
		20 \cdot 0.38 +
		10 \cdot 1 +
		15 \cdot 1
	}{
		20 + 10 + 15
	} = 0.71
\]

The normalized total score for this criterion is \emph{0.71}.

\paragraph{Ecosystem}
WebStorm provides full support for TypeScript and Node.js development, and the JetBrains plugin marketplace includes NestJS-specific plugins for code generation and navigation. \cite{bib:webstorm-ide}
The npm registry contains over 2.000.000 packages, a large number of which are compatible with NestJS applications. \cite{bib:npm-registry}

The scores for this criterion are shown in \cref{tab:multi-criteria-analysis-nestjs-ecosystem}.

\begin{htable}[
		caption = {NestJS ecosystem evaluation},
		label = {tab:multi-criteria-analysis-nestjs-ecosystem},
	]{
		colspec={lllX},
	}
	Sub-criterion & Weight & Function & Normalized score \\
	IDE plugin support & 25 & Binary (0 or 1) & 1 \\
	Number of libraries & 15 & $\bar{x} = \frac{N}{N_{\max}}$ & $\bar{x} = \frac{2000}{2000} = 1$ \\
\end{htable}

\[
	S_{\text{ecosystem}} =
	\frac{
		25 \cdot 1 +
		15 \cdot 1
	}{
		25 + 15
	} = 1.00
\]

The normalized total score for this criterion is \emph{1.00}.

\paragraph{Quality Assurance}
NestJS integrates with Jest out of the box and provides utilities for unit testing individual components and end-to-end testing the full application. \cite{bib:nestjs-testing}
The TypeScript ecosystem provides ESLint with TypeScript-specific rules, and NestJS projects are scaffolded with ESLint configuration by default. \cite{bib:nestjs-linting}
The last minor release of NestJS was within the last six months. \cite{bib:nestjs-github}

The scores for this criterion are shown in \cref{tab:multi-criteria-analysis-nestjs-quality-assurance}.

\begin{htable}[
		caption = {NestJS quality assurance evaluation},
		label = {tab:multi-criteria-analysis-nestjs-quality-assurance},
	]{
		colspec={lllX},
	}
	Sub-criterion & Weight & Function & Normalized score \\
	Testing & 15 & Binary (0 or 1) & 1 \\
	Static code analysis & 5 & Binary (0 or 1) & 1 \\
	Regular updates & 5 & Binary (0 or 1) & 1 \\
\end{htable}

\[
	S_{\text{quality assurance}} =
	\frac{
		15 \cdot 1 +
		5 \cdot 1 +
		5 \cdot 1
	}{
		15 + 5 + 5
	} = 1.00
\]

The normalized total score for this criterion is \emph{1.00}.
